\documentclass[10pt,a4paper]{amsart}
\usepackage[margin=19mm]{geometry}
\usepackage{amsmath,amssymb,mathtools}
\usepackage[scr=rsfs,cal=euler]{mathalfa}
\usepackage{xfrac}
\usepackage[unicode,hidelinks,bookmarks=false]{hyperref}
\newcommand{\ie}{i.e.\ }
\newcommand{\cf}{cf.\ }
\newcommand{\resp}{respectively\ }
\newcommand{\Subsection}{\subsection}
\providecommand{\nobreakdash}{\nobreak\hbox{-}}
\setlength{\emergencystretch}{2em}
\multlinegap0pt
\newtheorem{lemma}{Lemma}
\title{On invariant subrings of Orlik--Solomon and Varchenko--Gel'fand algebras in type A}
\author{Trevor Karn}
\date{Source: arXiv:2603.24884v1 (CC BY 4.0)}
\begin{document}
\maketitle

\noindent\emph{Source and licence.} On invariant subrings of Orlik--Solomon and Varchenko--Gel'fand algebras in type A. Version arXiv:2603.24884v1 (CC BY 4.0), distributed by arXiv under the Creative Commons Attribution 4.0 International licence, http://creativecommons.org/licenses/by/4.0/, as recorded in the arXiv licence metadata for this exact version and shown on the abstract page https://arxiv.org/abs/2603.24884v1. Full licence notice: \texttt{licenses/arxiv-2603.24884-CC-BY-4.0.txt}.\par\smallskip

\noindent\textit{Editorial scope.} Counted unit: the article's lemma lem:PartialIdentities (source0.tex lines 511--524). The article's complete original proof of it is delivered with it (source0.tex lines 525--575, one \texttt{proof} environment of 51 lines). No mathematics is written, repaired or re-derived: the statement and the proof are the article's own lines, in the article's own notation, and no retained line is substituted or re-flowed.  No further source-local context is needed: every cross-reference of every retained line resolves inside the two delivered spans.  Nothing was omitted from any retained span, so the package carries no omission record.  No retained line contains a citation, so the wrapper carries no bibliography and no bibliography entry.  No retained line contains a figure environment, an image-inclusion command or drawing code, so no image is delivered and the package declares no asset dependency.  Editorial formatting only, in addition: an AMS \texttt{amsart} preamble that restates the article's own declarations and macros used by the retained lines, namely the environments \texttt{lemma} on separate counters, together with the standard packages those lines need.  The retained environments print in this document as Lemma 1 in order of appearance, whereas the article prints the same items with the numbering of its own full document.  The complete native source of the article as distributed in the arXiv e-print for this exact version is bundled unchanged as \texttt{sources/197174/source0.tex}.
\begin{lemma}\label{lem:PartialIdentities}

    \begin{enumerate}
        \item $\partial(a) = \binom{n}{2}$
        \item $\partial(m) = n$
        \item $\partial(am) = -na + \binom{n}{2}m$
        \item $\partial(c^d) = -2dac^{d-1} + d(n-1)mc^{d-1}$
        \item $\partial(amc^d) = -nac^d + \binom{n}{2}mc^d$
        \item $\partial(g) = \begin{cases}
            0 & n~\text{even}\\
            a - (p-1)m& n~\text{odd}
        \end{cases}$
    \end{enumerate}
\end{lemma}

\begin{proof}
    Since $\partial(e_{ij})=1$, we see that $\partial(a)$ (respectively $\partial(m)$) is an integer counting the number of monomials in $a$ (respectively $m$) which is $\binom{n}{2}$ (respectively $n$).
    Then
    \begin{align*}
        \partial(am)
        &= \partial(a)m - a \partial(m)\\
        &= \binom{n}{2}m - na
    \end{align*}

    The Leibniz rule implies the power rule that $\partial(c^d) = dc^{d-1}\partial(c)$.
    Thus, we compute
    \begin{align*}
        \partial(c) &= \partial\left( \sum_{1 \leq i < j \leq n} e_{ij}e_{i,n+1} + e_{ij}e_{j,n+1}\right) \\
        &=\sum_{1 \leq i < j \leq n} e_{i,n+1} - e_{ij} + e_{j,n+1} - e_{ij}\\
        &=(n-1) m - 2a.
    \end{align*}
    So
    \[\partial(c^d) = dc^{d-1}\left((n-1)m - 2a\right).\]
    Next,
    \begin{align*}
        \partial(amc^d)
        &= \partial(am)c^d + am \partial(c^d)\\
        &= \left( -na + \binom{n}{2}m \right) c^d + am \left(-2dac^{d-1} + d(n-1)mc^{d-1}\right)\\
        &= \left( -na + \binom{n}{2} m \right) c^d.
    \end{align*}
    Lastly,
    \begin{align*}
        \partial(g)
        &= \partial(am - pc)\\
        &= (-na + \binom{n}{2}m) - p \partial(c^1)\\
        &= (-na + \binom{n}{2}m) - p \left(-2a + (n-1)m \right)\\
        &= \left(-n +2p\right) a + \left(\binom{n}{2} -(n-1)p\right) m.
    \end{align*}
    When $n$ is even, $p = \lfloor\frac{n+1}{2} \rfloor = \lfloor\frac{n}{2}\rfloor = \frac{n}{2}$.
    When $n$ is odd, $p = \frac{n+1}{2}$.
    Thus,
    \begin{align*}
        \partial(g)
        &=
        \begin{cases}
           \left (-n+ 2 \frac{n}{2} \right)a + \left( \frac{n(n-1)}{2} - (n-1)\frac{n}{2}\right)m & n~\text{even}\\
           \left (-n+ 2 \frac{n+1}{2} \right)a + \left( \frac{n(n-1)}{2} - (n-1)\frac{n+1}{2}\right)m  & n~\text{odd}
        \end{cases}\\
        &=
        \begin{cases}
            0 & n~\text{even}\\
            a - (p-1)m& n~\text{odd}
        \end{cases}
        \qedhere
    \end{align*}
\end{proof}

\end{document}
